\documentclass[11pt]{article}
\usepackage[a4paper,margin=22mm]{geometry}
\usepackage{fontspec}
\setmainfont{DejaVu Serif}
\setsansfont{DejaVu Sans}
\usepackage{microtype}
\usepackage{xcolor}
\usepackage{hyperref}
\usepackage{enumitem}
\usepackage{parskip}
\definecolor{Accent}{HTML}{971D32}
\definecolor{Muted}{HTML}{666666}
\hypersetup{colorlinks=true,urlcolor=Accent,linkcolor=Accent}
\setlist{leftmargin=1.6em,itemsep=0.3em,topsep=0.35em}
\setcounter{tocdepth}{2}
\renewcommand{\contentsname}{Contents}
\newcommand{\sectionrule}{\par\vspace{-0.5em}\noindent\textcolor{Accent}{\rule{\linewidth}{0.6pt}}\par}
\begin{document}

\begin{center}
{\sffamily\bfseries\color{Accent} TOEFL WORD OF THE DAY\par}
\vspace{8mm}
{\fontsize{42}{48}\selectfont\bfseries turmoil\par}
\vspace{2mm}
{\Large\sffamily\color{Accent} /ˈtɝːmɔɪl/\par}
\vspace{2mm}
{\small\sffamily Local audio phrase: turmoil\par}
{\small\sffamily Audio file: audios/turmoil.mp3\par}
\vspace{2mm}
{\large\sffamily\color{Muted} uncountable noun\par}
\vspace{5mm}
{\sffamily serious confusion or disorder\par}
\vspace{6mm}
{\large Turmoil describes a period or condition of serious confusion, uncertainty, disorder, or conflict.\par}
\vspace{6mm}
\begin{minipage}{0.84\textwidth}
\itshape “The sudden leadership change threw the organization into turmoil.”
\end{minipage}\par
\vspace{10mm}
\textbf{Date:} August 28th 2026\quad\textbullet\quad \textbf{Level:} B1--mid-B2 TOEFL
\vfill
Created and compiled by \textbf{Amir Kooshky}\par
\vspace{2mm}
\href{https://t.me/KooshkyTOEFL}{Telegram Channel}\quad\textbullet\quad
\href{https://www.instagram.com/kooshkytoefl}{Instagram}\quad\textbullet\quad
\href{https://t.me/KooshkyTOEFL\_pv}{Personal Telegram}
\end{center}
\newpage
\tableofcontents
\newpage

\section{Meanings and usage}
\sectionrule
\subsection{Meaning 1: Serious confusion or disorder}
\textbf{Part of speech:} uncountable noun

\textbf{IPA:} /ˈtɝːmɔɪl/

\textbf{Pronunciation phrase:} turmoil

\textbf{Local audio file:} audios/turmoil.mp3

\textbf{Register:} neutral to formal; common in news, history, politics, economics, business, and academic writing

\textbf{Definition:} a state of great confusion, uncertainty, or disorder, especially during a period of serious change or conflict

\textbf{In simpler words:} serious confusion and disorder

\textbf{Typical pattern:} in turmoil; throw or plunge something into turmoil; political or economic turmoil

\subsubsection{Natural examples}
\begin{enumerate}
  \item The country experienced years of political turmoil after the government collapsed.
  \item Financial turmoil caused investors to become much more cautious.
  \item The sudden leadership change threw the organization into turmoil.
  \item The region remained in turmoil following months of violent conflict.
  \item Economic turmoil led many businesses to postpone major investments.
  \item The unexpected announcement created considerable turmoil within the university.
  \item The industry went through a period of turmoil as new technologies disrupted traditional business models.
  \item Social turmoil often increases when people lose confidence in public institutions.
\end{enumerate}

\subsubsection{Nuance and implication}
\textbf{Central nuance:} Turmoil suggests more than ordinary confusion. It usually describes a serious and unstable situation in which normal order has been disrupted.

\textbf{Usual implication:} The situation is often changing rapidly, difficult to control, and stressful for the people or institutions involved.

\textbf{Compared with stability:} Stability means conditions are relatively steady and predictable. Turmoil means conditions are highly unsettled, confused, or disorderly.

\subsubsection{Grammar notes}
\textbf{How to use it:} Turmoil is normally used as an uncountable noun, so we usually do not say a turmoil or turmoils.

\textbf{What can follow it:} It often appears after in, into, amid, or during. It is also frequently modified by words such as political, economic, financial, social, or emotional.

\textbf{Common preposition:} Use in turmoil for a state of disorder and into turmoil after verbs such as throw or plunge.

\textbf{Noun use:} Turmoil is usually uncountable. Say political turmoil, a period of turmoil, or the country was in turmoil rather than a turmoil.

\textbf{Position in a sentence:} It commonly appears after be or remain in, and after verbs such as cause, create, throw, or plunge.

\subsubsection{Useful patterns}
\textbf{Pattern:} be in turmoil

\textbf{Use:} Use this when a country, organization, market, or other system is experiencing serious disorder or uncertainty.

\textbf{Example:} The company was in turmoil after several senior executives resigned.

\textbf{Pattern:} throw + something + into turmoil

\textbf{Use:} Use this when an event suddenly creates serious confusion or instability.

\textbf{Example:} The unexpected court ruling threw the industry into turmoil.

\textbf{Pattern:} plunge + something + into turmoil

\textbf{Use:} Use this for a sudden and serious move into disorder or instability.

\textbf{Example:} The political crisis plunged the country into turmoil.

\textbf{Pattern:} political or economic turmoil

\textbf{Use:} Use these common combinations to describe instability in government or the economy.

\textbf{Example:} Years of economic turmoil reduced public confidence in financial institutions.

\textbf{Pattern:} amid or during turmoil

\textbf{Use:} Use this when something happens while serious disorder or uncertainty is continuing.

\textbf{Example:} Several companies failed amid the financial turmoil.

\subsubsection{Common wrong usage}
\paragraph{Correction 1}
\textbf{Wrong:} The country experienced a turmoil after the election.

\textbf{Correct:} The country experienced turmoil after the election.

\textbf{Why:} Turmoil is normally uncountable, so do not usually use a before it.

\paragraph{Correction 2}
\textbf{Wrong:} The company was on turmoil after the director resigned.

\textbf{Correct:} The company was in turmoil after the director resigned.

\textbf{Why:} The standard expression is in turmoil.

\paragraph{Correction 3}
\textbf{Wrong:} The announcement threw the company in turmoil.

\textbf{Correct:} The announcement threw the company into turmoil.

\textbf{Why:} Use into turmoil after throw or plunge.

\paragraph{Correction 4}
\textbf{Wrong:} There were several turmoils in the economy.

\textbf{Correct:} There were several periods of turmoil in the economy.

\textbf{Why:} Turmoil is normally uncountable, so use periods of turmoil rather than turmoils.

\section{Word family}
\sectionrule
\subsection{tumultuous}
\textbf{Part of speech:} adjective

\textbf{IPA:} /tuːˈmʌltʃuəs/

\textbf{Definition:} full of confusion, disorder, strong emotion, or rapid change

\textbf{Relationship to the headword:} Tumultuous is not a direct morphological form of turmoil, but it is a closely related academic adjective often used for situations characterized by turmoil.

\textbf{Grammar pattern:} tumultuous + period, year, history, relationship, or change

\textbf{Usage note:} Use tumultuous to describe a period or situation; use turmoil as the noun for the state of disorder itself.

\subsubsection{Examples}
\begin{enumerate}
  \item The country went through a tumultuous period of political change.
  \item The company had a tumultuous year marked by resignations and financial losses.
  \item The reform was introduced during a tumultuous period in the nation's history.
\end{enumerate}

\section{Common collocations}
\sectionrule
\subsection{political turmoil}
\textbf{Applies to:} Serious confusion or disorder

\textbf{Meaning:} serious instability or disorder in government or politics

\textbf{Example:} The country entered a period of political turmoil after the disputed election.

\subsection{economic turmoil}
\textbf{Applies to:} Serious confusion or disorder

\textbf{Meaning:} serious instability in the economy

\textbf{Example:} The recession created economic turmoil across the region.

\subsection{financial turmoil}
\textbf{Applies to:} Serious confusion or disorder

\textbf{Meaning:} serious instability in financial markets or institutions

\textbf{Example:} Financial turmoil caused several banks to restrict lending.

\subsection{social turmoil}
\textbf{Applies to:} Serious confusion or disorder

\textbf{Meaning:} serious disorder or tension within society

\textbf{Example:} Rapid economic change contributed to social turmoil.

\subsection{in turmoil}
\textbf{Applies to:} Serious confusion or disorder

\textbf{Meaning:} in a state of serious confusion or instability

\textbf{Pattern:} be or remain in turmoil

\textbf{Example:} The organization was in turmoil after the sudden resignation of its president.

\subsection{throw into turmoil}
\textbf{Applies to:} Serious confusion or disorder

\textbf{Meaning:} suddenly cause serious confusion or instability

\textbf{Pattern:} throw + person, organization, country, or system + into turmoil

\textbf{Example:} The announcement threw the entire industry into turmoil.

\subsection{plunge into turmoil}
\textbf{Applies to:} Serious confusion or disorder

\textbf{Meaning:} cause something to enter a state of serious disorder

\textbf{Pattern:} plunge + something + into turmoil

\textbf{Example:} The collapse of the coalition plunged the government into turmoil.

\subsection{period of turmoil}
\textbf{Applies to:} Serious confusion or disorder

\textbf{Meaning:} a period during which serious confusion or instability exists

\textbf{Example:} The institution survived a long period of turmoil.

\subsection{amid turmoil}
\textbf{Applies to:} Serious confusion or disorder

\textbf{Meaning:} while serious disorder or uncertainty is occurring

\textbf{Example:} The company continued operating amid considerable financial turmoil.

\section{Synonyms and antonyms}
\sectionrule
\subsection{Close synonyms}
\subsubsection{chaos}
\textbf{Part of speech:} uncountable noun

\textbf{Applies to:} Serious confusion or disorder

\textbf{Shared meaning:} Both can describe serious disorder and confusion.

\textbf{Difference:} Chaos suggests an extreme lack of order or organization. Turmoil often includes instability, conflict, and uncertainty and does not always mean complete disorder.

\textbf{Example:} The sudden evacuation created chaos at the airport.

\subsubsection{upheaval}
\textbf{Part of speech:} countable or uncountable noun

\textbf{Applies to:} Serious confusion or disorder

\textbf{Shared meaning:} Both often describe unstable periods of major social, political, or economic change.

\textbf{Difference:} Upheaval emphasizes a major change or disruption to an existing system. Turmoil emphasizes the confusion, uncertainty, and instability that may accompany such change.

\textbf{Example:} Industrialization caused major social upheaval in many communities.

\subsubsection{disorder}
\textbf{Part of speech:} uncountable noun

\textbf{Applies to:} Serious confusion or disorder

\textbf{Shared meaning:} Both can describe a situation in which normal order has broken down.

\textbf{Difference:} Disorder is broader and simply means lack of order. Turmoil usually suggests a more intense, unstable, and emotionally or politically charged situation.

\textbf{Example:} The collapse of local authority led to widespread disorder.

\subsubsection{unrest}
\textbf{Part of speech:} uncountable noun

\textbf{Applies to:} Serious confusion or disorder

\textbf{Shared meaning:} Both can describe unstable social or political conditions.

\textbf{Difference:} Unrest usually refers specifically to social or political dissatisfaction that may lead to protests, strikes, or violence. Turmoil is broader and can also describe economic, institutional, or personal instability.

\textbf{Example:} Rising food prices contributed to widespread social unrest.

\subsection{Direct or useful antonyms}
\subsubsection{stability}
\textbf{Part of speech:} uncountable noun

\textbf{Applies to:} Serious confusion or disorder

\textbf{Opposition:} Stability is a condition in which things remain relatively steady and predictable, while turmoil involves serious uncertainty, disorder, or instability.

\textbf{Example:} The reforms were intended to restore economic stability.

\subsubsection{order}
\textbf{Part of speech:} uncountable noun

\textbf{Applies to:} Serious confusion or disorder

\textbf{Opposition:} Order means that people, systems, or activities are organized and controlled, while turmoil means normal order has been seriously disrupted.

\textbf{Example:} Authorities gradually restored order after the crisis.

\section{Words learners may confuse}
\sectionrule
\subsection{chaos}
\textbf{Part of speech:} uncountable noun

\textbf{Why learners confuse them:} Both describe serious disorder, but chaos is often more extreme.

\textbf{Key difference:} Chaos usually suggests almost complete disorder. Turmoil can describe serious instability and confusion even when some structure still remains.

\textbf{turmoil:} The company was in turmoil after its chief executive resigned.

\textbf{chaos:} The evacuation descended into chaos when the communication system failed.

\subsection{unrest}
\textbf{Part of speech:} uncountable noun

\textbf{Why learners confuse them:} Both are common in news and academic discussions of unstable situations.

\textbf{Key difference:} Unrest usually involves dissatisfaction among a group of people, often expressed through protests or conflict. Turmoil is broader and can describe political, economic, organizational, or emotional instability.

\textbf{turmoil:} The banking system was in turmoil after several institutions collapsed.

\textbf{unrest:} The government faced growing unrest among workers and students.

\subsection{upheaval}
\textbf{Part of speech:} countable or uncountable noun

\textbf{Why learners confuse them:} Both frequently appear in descriptions of major political, economic, and social change.

\textbf{Key difference:} Upheaval emphasizes a major disruption or change itself. Turmoil emphasizes the confused and unstable condition that results from or accompanies such disruption.

\textbf{turmoil:} The country remained in turmoil for months after the revolution.

\textbf{upheaval:} The revolution caused a major political upheaval.

\section{Etymology}
\sectionrule
\textbf{Origin:} Turmoil entered English several centuries ago, but its deeper origin is uncertain.

\textbf{Memory link:} Think of turmoil as a situation in which normal order has been badly disturbed and everything feels unstable.

\textbf{Accuracy note:} The early history of the word is not completely clear, so uncertain claims about a specific source language or original form are best avoided.

\section{Practice exercises}
\sectionrule
\subsection{Correct or incorrect?}
Decide whether each sentence is correct. Correct the inaccurate ones.

\begin{enumerate}
  \item The country experienced severe political turmoil after the election.
  \item The organization was on turmoil after its director resigned.
  \item The crisis threw the financial markets into turmoil.
  \item The country went through several turmoils during the decade.
  \item Economic turmoil can reduce investment and consumer confidence.
\end{enumerate}

\subsection{Collocation practice}
Complete each sentence with a natural collocation.

\begin{enumerate}
  \item The sudden resignation threw the entire organization \_\_\_\_\_ turmoil.
  \item The country remained \_\_\_\_\_ turmoil for several months after the crisis.
  \item The recession caused considerable economic \_\_\_\_\_.
  \item The institution survived a difficult period \_\_\_\_\_ turmoil.
\end{enumerate}

\subsection{Complete answer key}
\subsubsection{Correct or incorrect?}
\begin{enumerate}
  \item \textbf{Correct} \emph{Political turmoil is a common expression for serious instability in politics.}
  \item \textbf{Incorrect}. The organization was in turmoil after its director resigned. \emph{The standard expression is in turmoil.}
  \item \textbf{Correct} \emph{Throw something into turmoil is a common pattern meaning cause serious confusion or instability.}
  \item \textbf{Incorrect}. The country went through several periods of turmoil during the decade. \emph{Turmoil is normally uncountable, so use periods of turmoil rather than turmoils.}
  \item \textbf{Correct} \emph{Turmoil naturally describes serious instability in an economy.}
\end{enumerate}

\subsubsection{Collocation practice}
\begin{enumerate}
  \item \textbf{into.} The sudden resignation threw the entire organization into turmoil. \emph{Use throw something into turmoil.}
  \item \textbf{in.} The country remained in turmoil for several months after the crisis. \emph{In turmoil means in a state of serious confusion or instability.}
  \item \textbf{turmoil.} The recession caused considerable economic turmoil. \emph{Economic turmoil is a common collocation for serious instability in the economy.}
  \item \textbf{of.} The institution survived a difficult period of turmoil. \emph{A period of turmoil is a natural way to refer to a limited time of serious instability.}
\end{enumerate}

\vfill
\begin{center}
\small\color{Muted} Created and compiled by \textbf{Amir Kooshky}\\
\href{https://t.me/KooshkyTOEFL}{Telegram Channel}\quad\textbullet\quad
\href{https://www.instagram.com/kooshkytoefl}{Instagram}\quad\textbullet\quad
\href{https://t.me/KooshkyTOEFL\_pv}{Personal Telegram}
\end{center}
\end{document}
